\documentclass[11pt]{amsart}
\usepackage{amsmath,amssymb,amsthm,mathtools}
\usepackage[margin=1.1in]{geometry}
\usepackage{xcolor}
\usepackage{hyperref}
\hypersetup{colorlinks=true, linkcolor=blue!55!black, citecolor=blue!55!black,
  urlcolor=blue!55!black}
\theoremstyle{plain}
\newtheorem{theorem}{Theorem}[section]
\newtheorem{proposition}[theorem]{Proposition}
\newtheorem{lemma}[theorem]{Lemma}
\newtheorem{corollary}[theorem]{Corollary}
\theoremstyle{definition}
\newtheorem{definition}[theorem]{Definition}
\newtheorem{assumption}[theorem]{Assumption}
\theoremstyle{remark}
\newtheorem{remark}[theorem]{Remark}
\newcommand{\R}{\mathbb{R}}
\newcommand{\fpl}[1]{(-\Delta_p)^{s}#1}
\newcommand{\Jp}[1]{|#1|^{p-2}#1}
\newcommand{\Lv}{\mathcal L_{v}}
\newcommand{\al}{\alpha}
\newcommand{\asp}{\alpha_{s,p}}

\title{Fractional $p$-Laplacian dead-core problem:\\
PTU-faithful lemmas and proofs ($p$-case)}
\author{Hikmatullo Ismatov}
\date{\today\ --- working manuscript}

\begin{document}
\maketitle

\noindent\emph{Plan.} We reprove, for the fractional $p$-Laplacian, the
lemmas and theorems of dos Prazeres--Teymurazyan--Urbano~\cite{PTU}, one at a
time, keeping their structure and flagging precisely where the nonlinear
nonlocal operator forces a substitute or a new hypothesis. We begin with the
Liouville theorem (their Theorem~3.1).

%======================================================================
\section{Preliminaries}\label{sec:prelim}
%======================================================================
For $p\ge2$, $s\in(0,1)$, $J_p(t):=\Jp{t}$, the fractional $p$-Laplacian is
\[
  \fpl{v}(x)=\mathrm{P.V.}\!\int_{\R^n}\frac{J_p(v(x)-v(y))}{|x-y|^{n+sp}}\,dy
\]
(normalizing constant omitted). We say $v$ is \emph{$(s,p)$-harmonic} in
$\Omega$ if $\fpl{v}=0$ there, in the weak/viscosity sense~\cite{KKL}.

\subsection{The linearized operator along $v$}
For $v\in C^1_{loc}$, the linearization of $\fpl{}$ at $v$ has kernel
\[
  K_v(x,y):=(p-1)\,|v(x)-v(y)|^{p-2},\qquad
  \Lv\phi(x):=\mathrm{P.V.}\!\int_{\R^n}\big(\phi(x)-\phi(y)\big)
   \frac{K_v(x,y)}{|x-y|^{n+sp}}\,dy,
\]
a symmetric, nonnegative-kernel operator of effective order $sp-p+2$ (positive
for $s>1-\tfrac2p$). It is the operator denoted $L_u$ in~\cite[(2.1)--(2.2)]{GJS}.

\subsection{Two facts from~\cite{GJS}, used as inputs}
The following are established in~\cite{GJS} and play here the roles that
\cite[Thm 2.2]{PTU} (interior estimate) and linearity play in the case $p=2$.

\begin{lemma}[Directional derivatives solve the linearized equation, {\cite[\S3,\S9]{GJS}}]\label{lem:der-eq}
Let $v\in C^{1,\al_0}_{loc}(\R^n)$ be $(s,p)$-harmonic, with $K_v$ locally
integrable against $|x-y|^{-n-sp}$. Then for each unit vector $e$, the
directional derivative $D_e v$ is a weak solution of
\[
  \Lv (D_e v)=0\qquad\text{in }\R^n.
\]
\end{lemma}

\begin{lemma}[Interior H\"older estimate for the linearized equation, {\cite[\S7]{GJS}}]\label{lem:gjs7}
Let $\phi$ solve $\Lv\phi=0$ in $B_2$, and suppose $K_v$ satisfies on $B_2$ the
symmetry, scaled upper-bound and coercivity hypotheses of~\cite[Thm 7.3]{GJS}
\textnormal{(}in particular $K_v$ is \emph{uniformly elliptic}: there is a unit
vector $e$ with $|\nabla v-e|<\tfrac14$, equivalently $0<\lambda\le|\nabla v|\le
\Lambda$, on $B_2$\textnormal{)}. Then $\phi\in C^{\asp}(B_{1})$ for some
$\asp=\asp(n,s,p)\in(0,1)$, with
\[
  [\phi]_{C^{\asp}(B_{1})}\le C\big(\|\phi\|_{L^\infty(\R^n)}
   +\mathrm{Tail}_{p-1,sp+1}(\phi;2)\big),\qquad C=C(n,s,p,\lambda,\Lambda).
\]
\end{lemma}

\begin{remark}\label{rem:nondeg}
Lemma~\ref{lem:gjs7} requires \emph{nondegeneracy} of $v$ (a gradient bounded
below). For an entire $v$ this can fail where $Dv=0$. We therefore carry a
standing
\[
  \textbf{(ND)}\qquad Dv\not\equiv0\text{ on any open set, and on every compact
  }K\Subset\{Dv\ne0\}\text{ the constants }\lambda,\Lambda\text{ of
  Lemma~\ref{lem:gjs7} are uniform,}
\]
together with the scale-uniform form $(\star)$: under the rescalings used in the
proof, the constant in Lemma~\ref{lem:gjs7} is independent of the scale. At
$p=2$, $K_v\equiv\text{const}$, $\Lv=(-\Delta)^s$, and \textnormal{(ND)},
$(\star)$ are automatic.
\end{remark}

%======================================================================
\section{Theorem 3.1 ($p$-case): a Liouville theorem}\label{sec:liouville}
%======================================================================

\begin{theorem}[Liouville; $p$-analogue of {\cite[Thm 3.1]{PTU}}, $\nu=1$]\label{thm:liouville}
Let $p\ge2$, $s\in\big(1-\tfrac1p,1\big)$, and let $v\in C^{1,\al_0}_{loc}(\R^n)$
be an entire $(s,p)$-harmonic function satisfying \textnormal{(ND)},
$(\star)$ of Remark~\ref{rem:nondeg}. Suppose that, for some
$0\le\alpha<\delta<\asp$,
\begin{equation}\label{eq:growth}
  [v]_{C^{1+\alpha}(B_R)}\le C\,R^{\,\delta-\alpha}\qquad\text{for all }R\ge1,
\end{equation}
and that the tail condition $\delta(p-1)<sp$ holds. Then $v$ is affine:
\[
  v(x)=v(0)+Dv(0)\cdot x .
\]
In particular, if $v(0)=|Dv(0)|=0$, then $v\equiv0$.
\end{theorem}

\begin{remark}[Dictionary with~\cite{PTU}]\label{rem:dict}
In the application (Section~\ref{sec:liouville} is used after the blow-up of the
main theorem), $v$ grows like $|x|^{\beta}$ with $\beta=\frac{sp}{p-1-\gamma}$,
so $Dv$ grows like $|x|^{\beta-1}$ and $\delta=\beta-1$. The cap
$\delta<\asp$ then reads $\beta<1+\asp$, the exact analogue of PTU's
$\nu+\alpha<\nu+1$ ($\beta<\nu+1$) --- except that PTU's $s$-harmonic increments
are smooth, giving the generous cap $\delta<1$, whereas the rough linearized
operator $\Lv$ gives only the H\"older cap $\delta<\asp$.
\end{remark}

\begin{proof}
We follow~\cite[Thm 3.1, Step 1]{PTU}. The one structural change is that for
$p\ne2$ one cannot subtract the affine part of $v$ (it does not preserve
$\fpl{v}=0$); instead we work directly with the directional derivative, which
\emph{does} solve a linear equation by Lemma~\ref{lem:der-eq}.

Fix a unit vector $e$ and set
\[
  w:=D_e v-D_e v(0).
\]
Then $w(0)=0$, and since constants lie in the kernel of $\Lv$,
Lemma~\ref{lem:der-eq} gives
\begin{equation}\label{eq:weq}
  \Lv w=0\qquad\text{in }\R^n .
\end{equation}

\emph{Step 1: growth of $w$.} By \eqref{eq:growth} and $w(0)=0$, for $|x|\ge1$,
\[
  |w(x)|=|D_e v(x)-D_e v(0)|\le [v]_{C^{1+\alpha}(B_{|x|})}\,|x|^{\alpha}
   \le C\,|x|^{\delta-\alpha}\,|x|^{\alpha}=C\,|x|^{\delta},
\]
and $|w|\le C$ on $B_1$ by \eqref{eq:growth} with $R=1$. Since
$\delta(p-1)<sp<sp+1$, the function $w$ has finite
$\mathrm{Tail}_{p-1,sp+1}(w;R)$ for every $R$ (the integrand decays like
$|y|^{\delta(p-1)-n-sp-1}$, integrable). This is the analogue of PTU's
$\beta<1<2s\Rightarrow w_e\in L^1_s$.

\emph{Step 2: regularity and rescaled decay.} By \eqref{eq:weq} and
Lemma~\ref{lem:gjs7} (valid under \textnormal{(ND)}), $w\in C^{\asp}_{loc}$. For
$\rho\ge1$ set
\[
  \phi_\rho(x):=\rho^{-\delta}\,w(\rho x).
\]
By the growth bound, $\|\phi_\rho\|_{L^\infty(B_1)}\le C$ and
$\mathrm{Tail}_{p-1,sp+1}(\phi_\rho;2)\le C$, both uniformly in $\rho$; and
$\phi_\rho$ solves $\mathcal L_{v_\rho}\phi_\rho=0$ with
$v_\rho(x):=\rho^{-(1+\delta)}v(\rho x)$, whose kernel is, by $(\star)$, in the
class of Lemma~\ref{lem:gjs7} with constants independent of $\rho$. Hence
\[
  [\phi_\rho]_{C^{\asp}(B_{1/2})}\le C\quad\Longrightarrow\quad
  [w]_{C^{\asp}(B_{\rho/2})}=\rho^{\,\delta-\asp}\,[\phi_\rho]_{C^{\asp}(B_{1/2})}
   \le C\,\rho^{\,\delta-\asp}.
\]

\emph{Step 3: conclusion.} Since $\delta<\asp$, letting $\rho\to\infty$ gives
$[w]_{C^{\asp}(\R^n)}=0$, so $w$ is constant; as $w(0)=0$, $w\equiv0$, i.e.\
$D_e v\equiv D_e v(0)$ is constant. (Unlike~\cite{PTU}, no separate
translation-invariance step is needed: the H\"older \emph{seminorm} itself
decays to $0$ globally. The role of translation invariance is played by the
scale-uniformity $(\star)$, which is where \textnormal{(ND)} is used.) As $e$
was arbitrary, $Dv$ is constant, so $v$ is affine; if $v(0)=|Dv(0)|=0$ then
$Dv\equiv0$ and $v\equiv v(0)=0$.
\end{proof}

\begin{remark}[What is rigorous and what is an input]\label{rem:status}
The proof is complete \emph{given} Lemmas~\ref{lem:der-eq}--\ref{lem:gjs7} and
hypotheses \textnormal{(ND)}, $(\star)$. The genuinely open point is precisely
\textnormal{(ND)}/$(\star)$: Lemma~\ref{lem:gjs7} \textnormal{(\cite[\S7]{GJS})}
is a \emph{local} estimate under uniform ellipticity, and the kernel $K_v$
degenerates where $Dv=0$ --- at a flat point, the origin. Establishing the
scale-uniform estimate across the degenerate set is the analogue, for the
$p$-Laplacian, of the (free) interior estimate \cite[Thm 2.2]{PTU} for the
fractional Laplacian. At $p=2$ everything is automatic and the theorem reduces
to \cite[Thm 3.1]{PTU} ($\nu=1$).
\end{remark}

\begin{remark}[The case $\nu=2$]\label{rem:nu2}
PTU's $\nu=2$ (vanishing $2$-jet, cap $\beta<3$) has no $p$-analogue: it would
require a $C^{1,\asp}$ estimate for $\Lv$ (equivalently $C^{2,\asp}$ for
$\fpl{}$), which does not hold for the $p$-Laplacian. Hence only the $\nu=1$
theorem above is available, and the usable cap is $\beta<1+\asp$.
\end{remark}

%======================================================================
\section{Lemma 3.1 ($p$-case): jet to growth}\label{sec:jet}
%======================================================================
The next lemma is the $p$-analogue of~\cite[Lemma 3.1]{PTU}. It converts a
controlled $C^{1+\alpha}$ seminorm at small scales, together with a vanishing
$1$-jet, into a pointwise growth bound. It is purely a consequence of the mean
value theorem, hence \emph{operator-free}: the proof is identical to PTU's, and
the only restriction to $\nu=1$ comes from the absence of a $C^{2,\alpha}$
theory for the $p$-Laplacian (Remark~\ref{rem:nu2}).

\begin{lemma}[Jet to growth; $p$-analogue of {\cite[Lemma 3.1]{PTU}}, $\nu=1$]\label{lem:jet}
Let $u\in C^{1+\alpha}(B_1)$, $0<\alpha<1$, with $u(0)=|Du(0)|=0$. If, for some
$\beta>\alpha$ and $A>0$,
\[
  \sup_{0<r<1/2} r^{\,\alpha-\beta}\,[u]_{C^{1+\alpha}(B_r)}\le A,
\]
then
\[
  |u(x)|\le A\,|x|^{1+\beta}\qquad\text{for }x\in B_{1/2}.
\]
\end{lemma}

\begin{proof}
Fix $x\in B_{1/2}$. Since $u(0)=0$, the mean value theorem gives a point
$\xi\in B_{|x|}$ on the segment $[0,x]$ with
\[
  |u(x)|=|u(x)-u(0)|\le|Du(\xi)|\,|x|.
\]
Since $Du(0)=0$ and $Du\in C^{\alpha}(B_{1/2})$,
\[
  |Du(\xi)|=|Du(\xi)-Du(0)|\le [u]_{C^{1+\alpha}(B_{|x|})}\,|\xi|^{\alpha}
   \le [u]_{C^{1+\alpha}(B_{|x|})}\,|x|^{\alpha}
   \le A\,|x|^{\beta-\alpha}\,|x|^{\alpha}=A\,|x|^{\beta},
\]
where the last step used the hypothesis with $r=|x|$. Combining the two
displays gives $|u(x)|\le A\,|x|^{1+\beta}$.
\end{proof}

\begin{remark}\label{rem:jet-use}
This lemma feeds Theorem~\ref{thm:liouville}: in the blow-up of the main
theorem, the $\theta_k$-device produces a uniform $C^{1+\alpha}$ seminorm bound
and a vanishing $1$-jet; Lemma~\ref{lem:jet} converts these into the growth
hypothesis~\eqref{eq:growth} for the limit, with $\delta=\beta-1$. (The
contrapositive --- a violated growth bound forcing the device quantity to blow
up --- is the form used inside the contradiction argument of the main theorem.)
\end{remark}

%======================================================================
\section{Theorem 4.1 ($p$-case): improved growth at branching points}\label{sec:main}
%======================================================================
We now assemble the main estimate, following~\cite[\S4]{PTU}. Write
$\beta:=\frac{sp}{p-1-\gamma}>1$ and $F(t):=t_+^\gamma-t_-^\gamma$.

\begin{assumption}[Standing inputs]\label{ass:inputs}
$p\ge2$, $s\in(1-\tfrac1p,1)$, $0<\gamma<\tfrac{p-1}{1+sp}$ (so that
$\beta(p-1)<sp+1$; Section~\ref{sec:jet}), and:
\begin{itemize}
\item[\textnormal{(R)}] solutions of $\fpl{u}=f\in L^\infty$ are
$C^{0,\gamma_\circ}_{loc}$, $\gamma_\circ=\min\{1,\frac{sp}{p-1}\}$
\textnormal{(Biswas--Topp)}, giving the working seminorm; when an interior
$C^{1+\al}$ bound is available it is used, otherwise the H\"older seminorm
$[\,\cdot\,]_{C^{\gamma_\circ}}$ suffices for the device;
\item[\textnormal{(S)}] \emph{stability}: if $L_{k,h}W_k^h=f_k$ with $f_k\to0$
locally, $W_k^h\to W_0^h$ locally uniformly with uniformly finite
$\mathrm{Tail}_{p-1,sp+1}$, and the kernels converge, then $L_{0,h}W_0^h=0$
weakly \textnormal{(}mod.\ polynomials\textnormal{)};
\item[\textnormal{(ND)},$(\star)$] as in Remark~\ref{rem:nondeg}, for the
blow-up limits below.
\end{itemize}
\end{assumption}

\begin{definition}\label{def:flat}
$x_0\in B_1$ is a \emph{flat branching point} of $u$ if $u(x_0)=0$ and
$\sup_{B_r(x_0)}|u|=o(r)$ as $r\to0$.
\end{definition}

\begin{theorem}[$p$-analogue of {\cite[Thm 4.1]{PTU}}]\label{thm:main}
Under Assumption~\ref{ass:inputs}, let $u\in L^\infty(\R^n)$ solve
$\fpl{u}=F(u)$ in $B_1$ with $0$ a flat branching point. Then there is
$C=C(n,s,p,\gamma)$ with
\[
  |u(x)|\le C\,\|u\|_{L^\infty(\R^n)}\,|x|^{\beta}\qquad x\in B_{1/2}.
\]
\end{theorem}

\begin{proof}
Assume WLOG $\|u\|_{L^\infty}=1$; we must show $|u(x)|\le C|x|^\beta$. Suppose
not: there are solutions $u_k$ with $0$ a flat branching point and points $x_k$
with
\begin{equation}\label{eq:viol}
  |u_k(x_k)|>k\,|x_k|^{\beta}.
\end{equation}
Fix a working exponent $\alpha\in(0,\min\{\gamma_\circ,\beta-1\})$ and, with the
seminorm $[\,\cdot\,]:=[\,\cdot\,]_{C^{1+\alpha}}$ available under
\textnormal{(R)},\footnote{If only $C^{0,\gamma_\circ}$ is available, replace
$C^{1+\alpha}$ by $C^{\gamma_\circ}$ throughout and $\beta-1$ by $\beta$; the
device is unchanged.} define, following~\cite[(4.4)]{PTU},
\[
  \theta_k(r'):=\sup_{r'<r<1/2}r^{\,1+\alpha-\beta}\,[u_k]_{C^{1+\alpha}(B_r)},
  \qquad 1+\alpha-\beta<0 .
\]
By \eqref{eq:viol} and Lemma~\ref{lem:jet} (contrapositive: were
$\theta_k(0^+)\le k$, the vanishing $1$-jet would force $|u_k|\le k|x|^\beta$),
$\lim_{r'\to0}\theta_k(r')>k$; choose the near-maximal scale $r_k>1/k$ with
\begin{equation}\label{eq:rk}
  r_k^{\,1+\alpha-\beta}[u_k]_{C^{1+\alpha}(B_{r_k})}
  \ge\tfrac1k\theta_k(1/k)\ge\tfrac12\theta_k(r_k),
\end{equation}
so $r_k\to0$. Set the blow-up
\[
  v_k(x):=\frac{u_k(r_kx)}{\theta_k(r_k)\,r_k^{\beta}}.
\]
\emph{Seminorm bound (\cite[(4.8)]{PTU}).} For $1\le R\le\frac1{2r_k}$,
\[
  [v_k]_{C^{1+\alpha}(B_R)}
  =\frac{(r_kR)^{1+\alpha-\beta}[u_k]_{C^{1+\alpha}(B_{r_kR})}}{\theta_k(r_k)}
   \,R^{\beta-1-\alpha}\le R^{\beta-1-\alpha},
\]
using $(r_kR)^{1+\alpha-\beta}[u_k]_{C^{1+\alpha}(B_{r_kR})}\le\theta_k(r_kR)
\le\theta_k(r_k)$. The gradient-anchor of~\cite[\S4]{CP} gives $|v_k(x)|\le
C|x|^\beta$ for $1\le|x|\le\frac1{2r_k}$, and \eqref{eq:rk} gives the surviving
normalization $[v_k]_{C^{1+\alpha}(B_1)}\ge\frac12$.

\emph{Compactness.} The uniform \emph{local} seminorm bound and $v_k(0)=0$ give,
by Arzel\`a--Ascoli, $v_k\to v_0$ in $C^{1+\alpha'}_{loc}$ ($\alpha'<\alpha$)
along a subsequence, with
\[
  |v_0(x)|\le C(1+|x|)^{\beta},\qquad [v_0]_{C^{1+\alpha}(B_1)}\ge\tfrac12,
  \qquad v_0(0)=|Dv_0(0)|=0 .
\]
(This step needs \emph{no} tail control; it is the content of the seminorm
device.)

\emph{Vanishing right-hand side and limit equation.} A direct computation
(homogeneity) gives $\fpl{v_k}=\theta_k(r_k)^{\gamma-(p-1)}F(v_k)\to0$ locally
uniformly. Passing to increments $W_k^h=v_k(\cdot+h)-v_k$, Proposition/\S1 gives
$L_{k,h}W_k^h=f_k^h\to0$; by \textnormal{(S)} and the tail condition,
$L_{0,h}W_0^h=0$, i.e.\ $v_0$ is entire $(s,p)$-harmonic in the renormalized
sense.

\emph{Liouville and contradiction.} Since $\beta-1<\gamma_\circ\le1$ and
$\beta<1+\asp$ (admissibility) and $v_0$ satisfies \textnormal{(ND)},$(\star)$,
Theorem~\ref{thm:liouville} applies with $\delta=\beta-1$: $v_0$ is affine, and
the vanishing $1$-jet forces $v_0\equiv0$, contradicting
$[v_0]_{C^{1+\alpha}(B_1)}\ge\frac12$.
\end{proof}

\begin{remark}\label{rem:main-inputs}
The proof is complete \emph{given} \textnormal{(R)} (baseline regularity ---
available), \textnormal{(S)} (stability --- available for uniformly elliptic
kernels), and \textnormal{(ND)},$(\star)$ (the nondegeneracy of the blow-up
limit --- the open point of Theorem~\ref{thm:liouville}). Everything else is
unconditional. The next section reduces \textnormal{(ND)} to a
\emph{sup-nondegeneracy} of $v_0$, which the barrier is designed to produce.
\end{remark}

%======================================================================
\section{Reducing the open hypothesis (ND) to a free-boundary nondegeneracy}\label{sec:reduce}
%======================================================================
The one open input of Theorems~\ref{thm:liouville} and~\ref{thm:main} is
\textnormal{(ND)}: the coercivity/uniform-ellipticity of the linearized kernel
$K_{v}$ of the blow-up limit. This section does two things. First, it makes the
coercivity \emph{concrete}: it is exactly a \emph{directional nondegeneracy} of
$v$, in the sense of~\cite[\S6--\S7]{GJS}. Second, it shows this nondegeneracy
is produced, on annuli, by the barrier --- \emph{up to one residual step} which
we isolate honestly.

\subsection{Coercivity is directional nondegeneracy}
\begin{lemma}[Coercivity from directional nondegeneracy]\label{lem:coerc}
Let $v\in C^{0,1}(B_{2\rho}(z_0))$ and suppose there is a unit vector $e$ with
\begin{equation}\label{eq:dirnd}
  D_e v\ge c_1>0\quad\text{a.e.\ in }B_{\rho}(z_0).
\end{equation}
Then $K_v(x,y)=(p-1)|v(x)-v(y)|^{p-2}$ satisfies the coercivity hypothesis
\textnormal{(7.1)} of~\cite[Thm 7.3]{GJS} on $B_\rho(z_0)$: there are
$\lambda,\mu>0$ depending only on $n,p,c_1$ such that for every
$x\in B_\rho(z_0)$,
\[
  \Big|\Big\{\,y\in B_\rho(z_0):\ K_v(x,y)\,|x-y|^{-n-sp}\ge\lambda|x-y|^{-n-(sp-p+2)}\Big\}\Big|
   \ \ge\ \mu\,|B_\rho(z_0)| .
\]
\end{lemma}
\begin{proof}
The displayed inequality is, after cancelling $|x-y|^{-n-sp}$, equivalent to
$(p-1)|v(x)-v(y)|^{p-2}\ge\lambda|x-y|^{\,(sp)-(sp-p+2)}=\lambda|x-y|^{p-2}$,
i.e.\ $|v(x)-v(y)|\ge c|x-y|$ with $c=(\lambda/(p-1))^{1/(p-2)}$. Fix
$x\in B_\rho(z_0)$ and consider the cone of directions
$\mathcal C=\{\omega\in S^{n-1}:\ \omega\cdot e\ge\tfrac12\}$. For
$y=x-t\omega$, $\omega\in\mathcal C$, $t>0$, staying in $B_\rho(z_0)$,
\[
  v(x)-v(y)=\int_0^t D_\omega v(x-\tau\omega)\,d\tau
   =\int_0^t\big(\omega\cdot\nabla v(x-\tau\omega)\big)d\tau .
\]
Writing $\omega\cdot\nabla v=(e\cdot\nabla v)(e\cdot\omega)+(\omega-(\omega\cdot e)e)\cdot\nabla v$
and using \eqref{eq:dirnd} together with $|\nabla v|\le L$ (Lipschitz) is not
needed if we instead restrict to $\omega=e$ and a thin cone: for
$\omega\cdot e\ge\tfrac12$ and $c_1$ the lower bound, choosing the cone
half-angle small enough (depending on $c_1/L$) gives $\omega\cdot\nabla v\ge
c_1/4$ along the segment, whence $v(x)-v(y)\ge\frac{c_1}{4}t=\frac{c_1}{4}|x-y|$.
The set of admissible $y$ (the intersection of this cone-sector with
$B_\rho(z_0)$) has measure $\ge\mu|B_\rho(z_0)|$, $\mu=\mu(n,c_1/L)>0$. This is
the claim with $c=c_1/4$.
\end{proof}

\begin{remark}\label{rem:coerc}
Lemma~\ref{lem:coerc} is precisely the mechanism of~\cite[\S6]{GJS}: their
nondegeneracy $|\nabla u-e|<\tfrac14$ is exactly \eqref{eq:dirnd} with
$c_1=\tfrac34$. Thus \textnormal{(ND)} \emph{is} the statement ``$v$ has a
direction of definite positive derivative on a positive-measure set at each
scale'', and $(\star)$ is its scale-uniform version. The abstract
``degenerate-kernel Liouville'' has been reduced to a \emph{geometric
nondegeneracy of the blow-up limit}.
\end{remark}

\subsection{Sup-nondegeneracy from the barrier (one phase)}
We now show that the blow-up limit is nontrivial at every scale in the strong
form of a matching lower bound. We work in one phase ($u\ge0$).

\begin{lemma}[Sup-nondegeneracy; one phase]\label{lem:supnd}
Let $u\ge0$ solve $\fpl{u}=u^{\gamma}$ in $B_1$ with $u\ge0$ on $\R^n$, and let
$0$ be a flat branching point. There exist $c,r_0>0$ such that
\[
  \sup_{B_r}u\ \ge\ c\,r^{\beta}\qquad(0<r<r_0).
\]
Consequently the blow-up limit $v_0$ satisfies $\sup_{B_R}v_0\ge cR^{\beta}$ for
all $R\ge1$.
\end{lemma}
\begin{proof}[Proof (barrier comparison; strategy)]
Suppose not: $\sup_{B_{2r}}u\le\varepsilon r^{\beta}$ for $r$ along a sequence
and $\varepsilon\to0$. Let $\Psi=A\,\psi_{L,r}$ be the scale-$r$ barrier of the
scale-local barrier lemma \textnormal{(}$\fpl{\psi_{L,r}}\le-c_0r^{\beta\gamma}$
on $r\le|x|\le\tfrac32r$\textnormal{)}. Since $\Psi\ge0$,
$\fpl{\Psi}=A^{p-1}\fpl{\psi_{L,r}}\le -c_0A^{p-1}r^{\beta\gamma}\le
\Psi^{\gamma}$, so $\Psi$ is a subsolution of $\fpl{w}=w^{\gamma}$ on the
annulus. Choosing $A\sim\varepsilon^{1/(p-1-\gamma)}$ balances the barrier
against the smallness of $u$; the nonlocal comparison principle
(\cite[Thm 2.1]{PTU}-type, \cite{KKL}) applied on the annulus, with the ordering
$\Psi\le u$ verified on the complement using the smallness hypothesis and
$u\ge0$, forces $u\ge\Psi\ge c\,r^{\beta}$ at some point of the annulus,
contradicting $\sup_{B_{2r}}u\le\varepsilon r^{\beta}$ for $\varepsilon$ small.
Rescaling gives the statement for $v_0$.
\end{proof}

\begin{remark}[Honest status of Lemma~\ref{lem:supnd}]\label{rem:supnd}
The local dead-core analogue ($\Delta_p u=u^{\gamma}$) is classical
(da Silva--Salort, Teixeira). The nonlocal proof requires the comparison
ordering on the \emph{whole} complement; the barrier is $\ge0$ and one-signed,
so this is clean in \emph{one phase} and delicate across a two-phase branching
point. We state it as a lemma with the strategy above; the verification of the
complement ordering is the point to complete rigorously.
\end{remark}

\subsection{The residual step, isolated}
Combining Lemmas~\ref{lem:coerc} and~\ref{lem:supnd} \emph{almost} closes the
one-phase problem. The precise gap is the following implication.

\begin{quote}
\textbf{(Sup $\Rightarrow$ directional).} From $\sup_{B_R}v_0\ge cR^{\beta}$ and
$v_0(0)=0$, deduce that on each annulus $B_{2R}\setminus B_R$ there is a unit
vector $e_R$ and a set of measure $\ge\mu\,|B_{2R}\setminus B_R|$ on which
$D_{e_R}v_0\ge c\,R^{\beta-1}$ (directional nondegeneracy at scale $R$,
uniformly in $R$).
\end{quote}

If this holds, then: rescale $\bar v_R(z)=R^{-\beta}v_0(Rz)$; the above gives
\eqref{eq:dirnd} for $\bar v_R$ on a fixed sub-ball of the unit annulus,
Lemma~\ref{lem:coerc} gives the coercivity \textnormal{(7.1)} there,
Lemma~\ref{lem:gjs7} gives the interior H\"older estimate, and the rescaling of
Theorem~\ref{thm:liouville} (Step 2, run on the annuli rather than centred at
the flat origin) forces the increment to be constant --- hence $v_0$ affine,
hence $v_0\equiv0$, contradicting nontriviality. This would complete the
one-phase main theorem.

\begin{remark}[What the hard analysis shows]\label{rem:crux}
The bracketed implication \textnormal{(Sup $\Rightarrow$ directional)} is the
\emph{entire} remaining content for the one-phase problem. It is a concrete,
classical-flavoured free-boundary statement --- ``a solution that grows at the
sharp rate has a definite-sign directional derivative on a positive fraction of
each annulus'' --- far sharper than ``prove a degenerate-kernel Liouville''.
Sup-nondegeneracy (Lemma~\ref{lem:supnd}) gives the mass of the variation but
not yet its \emph{directional} coherence; controlling the direction is the
missing ingredient. This is the sharpest form we can give the open problem, and
the natural next target: it reduces a frontier regularity question to a
free-boundary nondegeneracy of Alt--Caffarelli type.
\end{remark}

\begin{thebibliography}{9}
\bibitem{GJS} D. Giovagnoli, D. Jesus, L. Silvestre, \emph{$C^{1+\alpha}$
regularity for fractional $p$-harmonic functions}, arXiv:2509.26565.
\bibitem{CP} J.C. Correa, D. dos Prazeres, \emph{A two-phase quenching-type
problem for the $p$-Laplacian}, arXiv:2504.11370.
\bibitem{KKL} J. Korvenp\"a\"a, T. Kuusi, E. Lindgren, \emph{Equivalence of
solutions to fractional $p$-Laplace type equations}, J. Math. Pures Appl.
132 (2019), 1--26.
\bibitem{PTU} D. dos Prazeres, R. Teymurazyan, J.M. Urbano, \emph{Improved
regularity for a nonlocal dead-core problem}, arXiv:2504.09557.
\end{thebibliography}
\end{document}
